% Generated aggregate supplement; include inside an existing Chinese LaTeX document.
\section{补充实验：预测边界与数值一致性}
所有新增检验均为既有数据上的复现和事后诊断，不构成新的未触碰测试。
\subsection{更强可观测基线}
同一观测编码和42日地标下，HMM与各基线条件对数似然差如下（nat/观测）。正数有利于HMM。
\begin{center}\begin{tabular}{lrrr}\hline 基线 & 差值 & 95\%下界 & 95\%上界\\\hline
IID & 0.142073 & 0.135053 & 0.149268 \\
observable-MC-order-1 & 0.022559 & 0.020116 & 0.025006 \\
observable-MC-order-2 & -0.008042 & -0.010232 & -0.005846 \\
observable-MC-order-3 & -0.024069 & -0.026859 & -0.021402 \\
\hline\end{tabular}\end{center}
二阶与三阶可观测链优于当前HMM，因此本文不主张隐状态模型具有普遍预测优势。
\subsection{OULAD配对区间}
按学生进行2000次配对重采样。第28日两类结局的AUC增益区间均跨0；Brier差为负表示改善。该区间条件于已拟合模型，不包含选模不确定性。
\begin{center}\begin{tabular}{llrrr}\hline 结局 & 指标 & 差值 & 下界 & 上界\\\hline
终局不利结果 & AUC & 0.002961 & -0.000792 & 0.006746 \\
终局不利结果 & AP & 0.003796 & 0.000049 & 0.007547 \\
终局不利结果 & Brier & -0.001021 & -0.001834 & -0.000194 \\
未来28日注销 & AUC & 0.016479 & -0.000775 & 0.033584 \\
未来28日注销 & AP & 0.000448 & -0.012540 & 0.011820 \\
未来28日注销 & Brier & -0.000082 & -0.000153 & -0.000011 \\
\hline\end{tabular}\end{center}
\subsection{风险递推与条件情景}
在两个协变量设定、四个情景和五个连续无活动阈值下，共40组DP与MC比较，最大标准化数值差为1.961。每情景仿真100000路径。
采用共同随机数，以配对差$D$的$\operatorname{sd}(D)/\sqrt{N}$报告仿真误差，不将其解释为策略效应区间。
未来28日风险只计入完全处于未来窗口内的连续无活动事件；地标处历史连续计数重置。
原未来已观测行为路径属于回顾性条件情景；冻结地标历史是另一假设设定，二者均不是可识别的因果干预。
KU编号不能跨课程、学年关联同一学生，故无法验证真实学生不重叠。完成分数在2021存在明显校准漂移，不应被解释为可靠通过概率。
